\chapter{Resumen}

RedNova GT es una empresa guatemalteca ficticia que ofrece software de gestión comercial a pequeñas y medianas empresas. Su crecimiento exige separar lógicamente las funciones de Ventas, Tecnologías de la Información (TI), Data Center y acceso de visitantes. El presente documento propone trasladar ese ecosistema a Amazon Web Services (AWS), dentro de una nube privada virtual con direccionamiento IPv4 privado, controles de ruta y reglas de seguridad específicas.

La solución parte del bloque \code{10.10.0.0/24}. Mediante Variable Length Subnet Masking (VLSM) se reservan bloques diferentes según la demanda: \code{/26} para Visitas, \code{/26} para Ventas, \code{/27} para TI, \code{/28} para Data Center y \code{/28} para Administración. El último \code{/26} queda sin asignar para crecimiento. Esta distribución considera una particularidad esencial de AWS: en cada subred, las primeras cuatro direcciones y la última no pueden asignarse a recursos \cite{aws-subnet-sizing}. Por esa razón el cálculo operativo utiliza $2^{32-p}-5$ y no la fórmula tradicional $2^h-2$.

La arquitectura separa dos dominios de exposición. Visitas y Administración se asocian con una tabla de rutas pública; Ventas, TI y Data Center se asocian con una tabla privada que sólo conserva la ruta local de la VPC. Para la prueba de concepto, dos instancias pequeñas se ubican dentro de Ventas y reciben exclusivamente direcciones privadas. El acceso administrativo atraviesa un bastión temporal cuyo puerto SSH se limita a una sola dirección IPv4 \code{/32}. El tráfico ICMP entre las dos instancias se autoriza mediante un grupo de seguridad común.

La infraestructura se describe como código en Terraform y mantiene el cómputo deshabilitado de forma predeterminada para prevenir cargos involuntarios. El 22 de octubre de 2026 se ejecutó una ventana de prueba en \code{us-east-1}: Terraform creó 20 recursos, incluidas la VPC, cinco subredes, tres instancias EC2 y sus controles de red. Las instancias privadas de Ventas intercambiaron cuatro paquetes ICMP con 0\% de pérdida, mientras que el mismo intento desde el bastión obtuvo 100\% de pérdida. La carpeta \code{evidence/} conserva los registros y capturas correspondientes.

\begin{infobox}[Resultado principal]
El diseño satisface las capacidades solicitadas, utiliza bloques canónicos sin solapamiento y conserva 64 direcciones para crecimiento. La subred \code{10.10.0.64/26} asignada a Ventas ofrece 59 direcciones utilizables en AWS. El prototipo confirmó en la nube las direcciones privadas \code{10.10.0.71} y \code{10.10.0.79}, junto con la política de aislamiento prevista.
\end{infobox}

El repositorio reúne diagramas, scripts, Terraform, LaTeX y evidencia de AWS. Con el despliegue y las pruebas P-01 a P-05 documentadas, el alcance técnico de la Fase 1 queda completo.
